% Generado por herramientas/ordenar_referencias.py: no editar a mano.
% Orden de primera cita (IEEE).
\begin{thebibliography}{35}
\bibitem{oms2022}
World Health Organization, \emph{WHO Recommendations for Care of the Preterm or Low-Birth-Weight Infant}. Geneva: WHO, 2022.

\bibitem{victora2021}
C.~G. Victora, P. Christian, L.~P. Vidaletti, G. Gatica-Domínguez, P. Menon, and R.~E. Black, ``Revisiting maternal and child undernutrition in low-income and middle-income countries: variable progress towards an unfinished agenda,'' \emph{The Lancet}, vol.~397, no.~10282, pp.~1382--1396, 2021.

\bibitem{fundacion2026}
Fundación Canguro, ``Proyecto final de la maestría en Inteligencia Artificial: Crecimiento y riesgos en el niño prematuro o de bajo peso al nacer (KMC-12mCA-Growing and Risks),'' propuesta y documentación institucional, 2026.

\bibitem{charpak2017}
N. Charpak, R. Tessier, J.~G. Ruiz \emph{et al.}, ``Twenty-year follow-up of kangaroo mother care versus traditional care,'' \emph{Pediatrics}, vol.~139, no.~1, e20162063, 2017.

\bibitem{charpak2023}
N. Charpak and A. Montealegre-Pomar, ``Follow-up of Kangaroo Mother Care programmes in the last 28 years: results from a cohort of 57\,154 low-birth-weight infants in Colombia,'' \emph{BMJ Global Health}, vol.~8, no.~5, 2023.

\bibitem{fenton2013}
T.~R. Fenton and J.~H. Kim, ``A systematic review and meta-analysis to revise the Fenton growth chart for preterm infants,'' \emph{BMC Pediatrics}, vol.~13, art.~59, 2013.

\bibitem{villar2015}
J. Villar, F. Giuliani, Z.~A. Bhutta \emph{et al.}, ``Postnatal growth standards for preterm infants: the Preterm Postnatal Follow-up Study of the INTERGROWTH-21st Project,'' \emph{The Lancet Global Health}, vol.~3, no.~11, pp.~e681--e691, 2015, doi: 10.1016/S2214-109X(15)00163-1.

\bibitem{oms2006}
World Health Organization, \emph{WHO Child Growth Standards: Length/Height-for-Age, Weight-for-Age, Weight-for-Length, Weight-for-Height and Body Mass Index-for-Age: Methods and Development}. Geneva: WHO, 2006.

\bibitem{ren2024}
Z. Ren, W. Gao, Q. Wang \emph{et al.}, ``Predictive role of NICU-related stress, postpartum depression trajectory and family coping on growth trajectory of moderate-to-late preterm infants: A longitudinal study,'' \emph{Journal of Advanced Nursing}, vol.~80, no.~8, pp.~3167--3178, 2024, doi: 10.1111/jan.16068.

\bibitem{nel2026}
S. Nel, U.~D. Feucht, T. Botha, M. Arashi, and F.~A.~M. Wenhold, ``First-year growth patterns of preterm infants receiving kangaroo mother care: associations with early life factors and 1-year anthropometry,'' \emph{European Journal of Clinical Nutrition}, vol.~80, pp.~62--72, 2026, doi: 10.1038/s41430-025-01662-6.

\bibitem{melo2016}
S.~E. Melo, O.~O. Melo, and N. Suarez, ``Modelo multinivel de tasa global de ganancia de peso en el programa madre canguro en Bogotá,'' \emph{Revista de Salud Pública}, vol.~18, no.~2, pp.~238--250, 2016, doi: 10.15446/rsap.v18n2.44661.

\bibitem{baker2023}
S. Baker and Y. Kandasamy, ``Machine learning for understanding and predicting neurodevelopmental outcomes in premature infants: a systematic review,'' \emph{Pediatric Research}, vol.~93, pp.~293--299, 2023.

\bibitem{demirci2024}
G.~M. Demirci \emph{et al.}, ``Predicting mental and psychomotor delay in very pre-term infants using machine learning,'' \emph{Pediatric Research}, vol.~95, pp.~668--678, 2024.

\bibitem{mlrisk2024}
``Machine learning risk stratification for high-risk infant follow-up of term and late preterm infants,'' \emph{Pediatric Research}, 2024, doi: 10.1038/s41390-024-03338-6.

\bibitem{wang2024}
L.-Y. Wang, L.-Y. Wang, M.-I. Sung \emph{et al.}, ``Implementing explainable machine learning models for practical prediction of early neonatal hypoglycemia,'' \emph{Diagnostics}, vol.~14, no.~14, art.~1571, 2024, doi: 10.3390/diagnostics14141571.

\bibitem{kahn2016}
M.~G. Kahn, T.~J. Callahan, J. Barnard \emph{et al.}, ``A harmonized data quality assessment terminology and framework for the secondary use of electronic health record data,'' \emph{eGEMs}, vol.~4, no.~1, art.~1244, 2016, doi: 10.13063/2327-9214.1244.

\bibitem{weiskopf2013}
N.~G. Weiskopf and C. Weng, ``Methods and dimensions of electronic health record data quality assessment: enabling reuse for clinical research,'' \emph{Journal of the American Medical Informatics Association}, vol.~20, no.~1, pp.~144--151, 2013.

\bibitem{chandola2009}
V. Chandola, A. Banerjee, and V. Kumar, ``Anomaly detection: A survey,'' \emph{ACM Computing Surveys}, vol.~41, no.~3, art.~15, 2009.

\bibitem{koenker1978}
R. Koenker and G. Bassett, ``Regression quantiles,'' \emph{Econometrica}, vol.~46, no.~1, pp.~33--50, 1978.

\bibitem{romano2019}
Y. Romano, E. Patterson, and E. Candès, ``Conformalized quantile regression,'' in \emph{Advances in Neural Information Processing Systems}, vol.~32, 2019.

\bibitem{angelopoulos2023}
A.~N. Angelopoulos and S. Bates, ``Conformal prediction: A gentle introduction,'' \emph{Foundations and Trends in Machine Learning}, vol.~16, no.~4, pp.~494--591, 2023.

\bibitem{vancalster2019}
B. Van Calster, D.~J. McLernon, M. van Smeden, L. Wynants, and E.~W. Steyerberg, ``Calibration: the Achilles heel of predictive analytics,'' \emph{BMC Medicine}, vol.~17, art.~230, 2019.

\bibitem{davis2017}
S.~E. Davis, T.~A. Lasko, G. Chen, and M.~E. Matheny, ``Calibration drift among regression and machine learning models for hospital mortality,'' in \emph{AMIA Annual Symposium Proceedings}, vol.~2017, pp.~625--634, 2017.

\bibitem{wong2021}
A. Wong, E. Otles, J.~P. Donnelly \emph{et al.}, ``External validation of a widely implemented proprietary sepsis prediction model in hospitalized patients,'' \emph{JAMA Internal Medicine}, vol.~181, no.~8, pp.~1065--1070, 2021, doi: 10.1001/jamainternmed.2021.2626.

\bibitem{page1954}
E.~S. Page, ``Continuous inspection schemes,'' \emph{Biometrika}, vol.~41, no.~1/2, pp.~100--115, 1954.

\bibitem{lundberg2020}
S.~M. Lundberg, G. Erion, H. Chen \emph{et al.}, ``From local explanations to global understanding with explainable AI for trees,'' \emph{Nature Machine Intelligence}, vol.~2, no.~1, pp.~56--67, 2020.

\bibitem{salih2024}
A.~M. Salih, G. Menegaz, T. Pillay, and E.~M. Boyle, ``Explainable artificial intelligence in paediatric: Challenges for the future,'' \emph{Health Science Reports}, vol.~7, no.~12, e70271, 2024, doi: 10.1002/hsr2.70271.

\bibitem{ghassemi2021}
M. Ghassemi, L. Oakden-Rayner, and A.~L. Beam, ``The false hope of current approaches to explainable artificial intelligence in health care,'' \emph{The Lancet Digital Health}, vol.~3, no.~11, pp.~e745--e750, 2021.

\bibitem{collins2024}
G.~S. Collins, K.~G.~M. Moons, P. Dhiman \emph{et al.}, ``TRIPOD+AI statement: updated guidance for reporting clinical prediction models that use regression or machine learning methods,'' \emph{BMJ}, vol.~385, e078378, 2024.

\bibitem{kreuzberger2023}
D. Kreuzberger, N. Kühl, and S. Hirschl, ``Machine learning operations (MLOps): Overview, definition, and architecture,'' \emph{IEEE Access}, vol.~11, pp.~31866--31879, 2023, doi: 10.1109/ACCESS.2023.3262138.

\bibitem{zaharia2018}
M. Zaharia, A. Chen, A. Davidson \emph{et al.}, ``Accelerating the machine learning lifecycle with MLflow,'' \emph{IEEE Data Engineering Bulletin}, vol.~41, no.~4, pp.~39--45, 2018.

\bibitem{designcouncil2007}
Design Council, \emph{Eleven Lessons: Managing Design in Eleven Global Companies}. London: Design Council, 2007.

\bibitem{brown2008}
T. Brown, ``Design thinking,'' \emph{Harvard Business Review}, vol.~86, no.~6, pp.~84--92, 2008.

\bibitem{ke2017}
G. Ke, Q. Meng, T. Finley \emph{et al.}, ``LightGBM: A highly efficient gradient boosting decision tree,'' in \emph{Advances in Neural Information Processing Systems}, vol.~30, pp.~3146--3154, 2017.

\bibitem{omsunicef2019}
World Health Organization and UNICEF, \emph{Recommendations for Data Collection, Analysis and Reporting on Anthropometric Indicators in Children under 5 Years Old}. Geneva: WHO, 2019.

\end{thebibliography}
